% Curs 10, pagina 2 -- continuare din curs10-pagina1-ro.tex.
% Fragment pentru \input; fara preambul, \clearpage sau sectiune de curs noua.

\subsection*{Curs 10 - Pagina 2}

\bigskip
\noindent\underline{Ex.} \textbf{3)} Fie $G$ grup arbitrar, $K,H\leq G$,
\quad $A=\{aH\mid a\in G\}$ -- clase de resturi stângi ale lui $H$ în $G$.

\bigskip
\noindent $K\times A\xrightarrow{\ \psi\ }A$.

\noindent $(k,aH)\overset{\psi}{\longmapsto}k\cdot aH=(ka)H$,
\quad $k\in K$, $aH\in A$.

\bigskip
\noindent $\psi$ -- acțiune prin translație la stânga pe mulțimea claselor
 de resturi ale lui $H$.

\bigskip
\noindent Dacă $A$ e o $G$-mulțime, îi putem asocia relația de
 echivalență ``$\sim_G$''.

\medskip
\noindent $x,y\in A$, \quad
$x\sim_G y\ \Longleftrightarrow\ (\exists)\,a\in G$ a.î. $ax=y$.

\bigskip
\noindent Dacă $x\in A$, \quad
$\widehat{x}\overset{\mathrm{def}}{=}\operatorname{orb}_G(x)
=Gx=\{ax\mid a\in G\}$ -- \underline{orbita lui $x$}.

\bigskip
\noindent\underline{Stabilizatorul în $G$ al lui $x$}:

\medskip
\noindent $\operatorname{Stab}_G(x)=\{a\in G\mid ax=x\}\leq G$.

\bigskip
\noindent Dacă $x\in A$, \quad
$\bigl|\operatorname{orb}_G(x)\bigr|
=[G:\operatorname{Stab}_G(x)]$.

\medskip
\noindent Dacă $|G|<\infty$, \quad
$\bigl|\operatorname{orb}_G(x)\bigr|\mid |G|$
(din teorema lui Lagrange).

\bigskip
\noindent\underline{Exp. (1)}:
\quad $\operatorname{Stab}_G(x)\overset{\mathrm{def}}{=}
\{a\in G\mid axa^{-1}=x\}
=\{a\in G\mid ax=xa\}$.

\medskip
\noindent În acest caz,
$\operatorname{Stab}_G(x)\overset{\mathrm{not.}}{=}C_G(x)$
-- \underline{centralizatorul în $G$ al lui $x$}.

\medskip
\noindent $\Longrightarrow\ \bigl|\operatorname{orb}_G(x)\bigr|
=[G:C_G(x)]$;
\quad $G$ comutativ $\Longrightarrow C_G(x)=G$.

\bigskip
\noindent\underline{Exp. (2)}: \quad $H\leq G$.

\medskip
\noindent $\operatorname{Stab}_G(H)\overset{\mathrm{def}}{=}
\{a\in G\mid a\cdot H=H\}
=\{a\in G\mid aHa^{-1}=H\}$.

\medskip
\noindent Notăm $N_G(H)$ -- \underline{normalizatorul în $G$ al lui $H$}.

\bigskip
\noindent\underline{Prop.}
\begin{enumerate}
\item[(1)] $H\subseteq N_G(H)$.
\item[(2)] $H\trianglelefteq N_G(H)$.
\item[(3)] $N_G(H)$ este cel mai mare subgrup al lui $G$
 pentru care $H$ este subgrup normal.
\end{enumerate}

\medskip
\noindent\underline{Dem.}

\noindent (1) \quad $h\in H$.

\noindent (iar $a\in H$).\footnote{Ultima paranteză din manuscris este greu de citit; ``iar $a\in H$'' este o transcriere provizorie. Demonstrația continuă pe pagina următoare.}
